\section{What the result establishes and what remains}

The exact law applies to every finite field, every feasible pair of frame dimensions, and every pair of restricted alternating forms. It gives probabilities of individual matrices, a complete conditional overlap law, and a necessary and sufficient support test. The minimum-dimension formula has a constructive witness. For fixed field and fixed nonempty frame dimensions, the discrepancy from uniform has a positive, explicitly computable leading coefficient at scale $q^{-2n}$.

This is a transfer of a finite method from the supplied GRH proof, not a consequence that needs RH as an assumption. The proof exposes a reusable distinction: the combined Gram matrix controls the geometric relations, while the intersections of its radical with the two coordinate blocks control whether each original frame remains injective after quotienting.

Three further questions remain separate. First, the constants in fixed-size asymptotics depend on the frame sizes, so those estimates alone do not settle every boundary regime with $a,b$ growing alongside $n$. Second, positive-type \emph{nonuniform} weights in the original source require representation-theoretic information beyond the orbit counts used here. Third, extension to symmetric or Hermitian forms needs the relevant isometry types, and in characteristic two quadratic data cannot be replaced by a symmetric matrix alone.

The contribution is a proved extension relative to the inspected research corpus, supported by a targeted literature comparison. External review could identify an equivalent enumeration in the broader finite-geometry literature; the classical ingredients have therefore been identified explicitly. This note does not assert resolution of any further Riemann hypothesis.
